\documentclass[11pt]{article}
\usepackage[T1]{fontenc}
\usepackage[utf8]{inputenc}
\usepackage[spanish,es-noshorthands]{babel}
\usepackage{lmodern}
\usepackage[margin=2.4cm]{geometry}
\usepackage{amsmath,amssymb}
\usepackage{booktabs}
\usepackage{xcolor}
\usepackage{enumitem}
\usepackage{parskip}
\usepackage[colorlinks=true,linkcolor=black,urlcolor=blue]{hyperref}

\newcommand{\campo}{E}

\title{\textbf{AstroBioSim — Modelo físico del entorno}\\[2pt]
\large Decisiones de diseño para la exposición y defensa final}
\author{Área de Física}
\date{Julio 2026}

\begin{document}
\maketitle

\begin{abstract}
\noindent El motor ambiental construye el campo físico que gobierna la
habitabilidad: temperatura $T$, radiación $R$ y actividad de agua $a_w$, alineadas
a la grilla del autómata. Este documento justifica las tres variables, los modelos
de campo de cada subsuelo (Tierra, Marte, Encelado), el reencuadre de la radiación
como UV, y los eventos estocásticos que perturban el ambiente. El criterio rector
es que \emph{una variable que no discrimina no entra al modelo}.
\end{abstract}

\section{El vector ambiental: por qué tres variables}

Cada celda lleva un vector $\campo_{i,j}=(T,\ R,\ a_w)$ con las tres capas alineadas
a la grilla $M\times N$. La especificación original tenía cuatro variables; la
\textbf{presión fue eliminada} (ADR-0008) porque en los análogos terrestres
permanece casi constante ($\sim 0{,}5$--$1$~atm), nunca cruza umbrales letales y
solo agrega una dimensión muerta. El principio se aplica de nuevo con la radiación
(§3): mantenemos únicamente lo que \emph{separa escenarios}.

\section{Modelos de campo por subsuelo}

Convención de grilla: la fila $0$ es la superficie y el eje de filas representa
\textbf{profundidad creciente} (salvo Encelado, que es un corte horizontal del
piso oceánico).

\paragraph{Tierra (control).} Subsuelo somero, térmicamente estable: la onda diurna
ya está amortiguada, de modo que el campo es uniforme a la media anual
($T\approx 19{,}8\,^\circ$C), con $R=0$ (el suelo bloquea el UV) y $a_w\approx
0{,}99$ (suelo a capacidad de campo). Es el escenario benigno de referencia. El
valor de $a_w$ es una corrección: el dataset traía la humedad relativa del
\emph{aire} (que oscilaba $0{,}16$--$0{,}93$), no la actividad de agua del suelo;
un subsuelo húmedo es estable y alto, no meteorológico.

\paragraph{Marte (análogo Atacama).} La onda térmica diurna se amortigua
\emph{exponencialmente} con la profundidad $z$ hacia una asíntota fría (modelo
estándar de \emph{skin depth} en regolito):
\begin{equation}
T(z)=T_{\infty}+(T_{\mathrm{sup}}-T_{\infty})\,e^{-z/\delta_T},\qquad
\mathrm{UV}(z)=\mathrm{UV}_{\mathrm{sup}}\,e^{-z/\delta_{\mathrm{UV}}},
\end{equation}
con $\delta_{\mathrm{UV}}\ll\delta_T$: la luz penetra el regolito muchísimo menos
que el calor. Esta separación de escalas es física central del proyecto (§4).

\paragraph{Encelado (océano subglacial).} Fondo frío ($T\approx 2{,}4\,^\circ$C)
con \emph{fumarolas} hidrotermales como fuentes puntuales cuyo pico de temperatura
decae radialmente con un núcleo gaussiano:
\begin{equation}
T(i,j)=T_{\text{océano}}+\sum_{f}\Delta T_f\,
\exp\!\left(-\frac{(i-i_f)^2+(j-j_f)^2}{2\sigma^2}\right),
\end{equation}
con $R=0$ (el hielo blinda por completo la radiación externa) y $a_w\approx 0{,}98$
(agua de mar). El radio $\sigma$ se fija como \emph{fracción de la grilla}, no en
celdas fijas: de lo contrario, en grillas pequeñas las fumarolas se solapan y el
pico depende de la resolución en vez de la física (corrección hallada en la
integración).

\section{Reencuadre de la radiación: UV, no dosis ionizante (ADR-0014)}

El dato crudo es irradiancia solar \emph{global} en W/m$^2$. Sin embargo, la
variable físicamente pertinente para la vida no es la insolación total (mayormente
visible e infrarroja, que no esteriliza) ni la \emph{dosis ionizante acumulada}.
Los números lo obligan:
\begin{itemize}[nosep]
  \item Superficie marciana: $0{,}21$~mGy/día $\approx 0{,}077$~Gy/año (RAD,
  Curiosity).
  \item \emph{D.\ radiodurans} tolera $5000$~Gy sin pérdida de viabilidad.
\end{itemize}
Acumular la dosis letal tomaría $\sim\!65\,000$~años: en la escala de la
simulación, la dosis ionizante \textbf{no discrimina} —el mismo defecto que la
presión. Lo que sí esteriliza de inmediato es el \textbf{UV}, y además se mide en
W/m$^2$, las mismas unidades. Por eso $R$ pasa a ser \textbf{irradiancia UV},
obtenida de la global por una fracción documentada,
$\mathrm{UV}=f_{\mathrm{UV}}\cdot R_{\text{global}}$ con $f_{\mathrm{UV}}=0{,}05$.

\paragraph{Validación independiente.} El modelo da $\mathrm{UV}_{\mathrm{sup}}
\approx 844\times 0{,}05 = 42{,}2$~W/m$^2$, que cae dentro del rango
\emph{publicado} de UV marciano (200--400~nm) de $42$--$55$~W/m$^2$. El factor no
estaba lejos por una vía distinta a la que lo fijó.

\section{La banda de profundidad habitable}

Combinando la atenuación rápida del UV con la lenta del calor, el eje de
profundidad adquiere significado biológico. Con los umbrales derivados de
literatura, en Marte emergen \emph{sin ajustar nada}:
\begin{center}\small
\begin{tabular}{@{}lll@{}}
\toprule
Profundidad & Estado & Causa \\
\midrule
Superficie (filas 0--2) & \textsf{MUERTA} & UV por encima de la fluencia letal \\
Somero (filas 3--4) & \textsf{LATENTE} & el UV ya no mata pero inhibe; falta agua \\
Profundo (filas 5$+$) & \textsf{LATENTE} & el UV deja de limitar; sigue faltando agua \\
\bottomrule
\end{tabular}
\end{center}
La ventana habitable marciana queda acotada \emph{arriba} por el UV y \emph{abajo}
por la disponibilidad de agua: una banda de profundidad, no un plano.

\section{Eventos estocásticos (Montecarlo)}

Perturbaciones aleatorias que actúan \emph{solo} sobre el campo (nunca sobre el
estado biológico), todas con el generador sembrado inyectado:
\begin{itemize}[nosep]
  \item \textbf{Micro-fisura marciana:} caída abrupta de $a_w$ en un radio de
  celdas (desecación por liberación de vapor).
  \item \textbf{Emisión hidrotermal (Encelado):} pico $\Delta T\sim\mathcal{N}(\mu,\sigma^2)$
  que se disipa radialmente con un núcleo gaussiano (difusión de calor).
  \item \textbf{Salmuera delicuescente (Marte):} contraparte simétrica de la
  micro-fisura que \emph{eleva} $a_w$ transitoriamente en un radio (ADR-0015).
\end{itemize}

\paragraph{Simetría necesaria.} Antes, todos los eventos \emph{degradaban} el
ambiente, sesgando cada corrida hacia la extinción por construcción. El evento de
salmuera es el que da a Marte algo que estudiar: sin él, la pregunta de
investigación no tiene mecanismo. Su frecuencia y magnitud son los parámetros del
barrido central del proyecto. Ya está implementado y verificado: sobre un refugio
transitorio (que sube $a_w$ hacia $0{,}90$--$0{,}98$ y decae con el tiempo),
\emph{D.\ radiodurans} recupera el crecimiento activo que el regolito en bulk le
niega.

\section{Dos modos de alimentar el campo (ADR-0017)}

El mismo motor corre bajo dos fuentes de campo intercambiables, sin duplicar el
bucle de simulación:
\begin{itemize}[nosep]
  \item \textbf{Sandbox:} campo homogéneo con $T$, $R$ y $a_w$ fijados a mano
  (sliders); explora el espacio de parámetros sin depender de datos.
  \item \textbf{Analógico:} la serie temporal real 2025. Cada dato diario se
  propaga con la \emph{estructura espacial} de estos modelos de subsuelo —el
  gradiente de profundidad en Marte, las fumarolas en Encelado— en vez de aplanar
  el campo. La decisión importa físicamente: un campo uniforme pondría el UV de
  superficie en todas las celdas y esterilizaría la grilla entera, borrando la
  banda de profundidad.
\end{itemize}

\section*{Deuda física abierta (honestidad para la defensa)}
\begin{itemize}[nosep]
  \item La $a_w$ de Marte usa el \emph{mínimo} diario real ($0{,}187$), una cota
  pesimista; conviene la media con su dispersión.
  \item La asíntota térmica $T_{\infty}=7{,}8\,^\circ$C es la media de los mínimos
  diarios, pero una onda térmica amortigua hacia la \emph{media anual} ($\approx
  22\,^\circ$C): revisar.
  \item El pico $\Delta T=25$ de la emisión hidrotermal, apilado sobre una fumarola
  existente, puede llevar $T$ a $\sim\!59\,^\circ$C. Puede ser correcto (el núcleo
  de una ventila es letal), pero hay que decidirlo.
\end{itemize}

\end{document}
